%% Some useful math commands

% Wide hat, bar, tilde
\newcommand\wdhat[1]{\hstretch{2}{\hat{\hstretch{.5}{#1}}}}
\newcommand\wdbar[1]{\hstretch{2}{\bar{\hstretch{.5}{#1}}}}
\newcommand\wdtilde[1]{\hstretch{2}{\tilde{\hstretch{.5}{#1}}}}

% \Set command
\providecommand\given{} % so it exists
\newcommand\SetSymbol[1][]{
   \nonscript\,#1\vert \allowbreak \nonscript\,\mathopen{}}
\DeclarePairedDelimiterX\Set[1]{\lbrace}{\rbrace}%
 { \renewcommand\given{\SetSymbol[\delimsize]} #1 }

% Absolute value, norm and inner product
\DeclarePairedDelimiterX\abs[1]{\lvert}{\rvert}{
\ifblank{#1}{\:\cdot\:}{#1}}
\DeclarePairedDelimiterX\norm[1]{\lVert}{\rVert}{
\ifblank{#1}{\:\cdot\:}{#1}}
\DeclarePairedDelimiterX\innerp[2]{\langle}{\rangle}{#1,#2}

%% Theorem environments and autoref
\newtheorem{thmdummy}{***}[chapter]
\newcommand{\mynewtheorem}[2]{%
  \newaliascnt{#1}{thmdummy}
  \newtheorem{#1}[#1]{#2}
  \aliascntresetthe{#1}
  \expandafter\def\csname #1autorefname\endcsname{#2}
}

\theoremstyle{plain}
    \mynewtheorem{theorem}{Theorem}
    \mynewtheorem{lemma}{Lemma}
    \mynewtheorem{corollary}{Corollary}
    \mynewtheorem{proposition}{Proposition}
    \newtheorem*{theorem*}{Theorem}
    \newtheorem*{proposition*}{Proposition}
    \newtheorem*{lemma*}{Lemma}

\theoremstyle{definition}
    \mynewtheorem{definition}{Definition}
    \mynewtheorem{example}{Example}
    \newtheorem*{example*}{Example}

\theoremstyle{remark}
    \mynewtheorem{remark}{Remark}
    \mynewtheorem{notation}{Notation}


% BibLaTeX bib library
\addbibresource{bib.bib}